% MPI_Scatter: root's send buffer is split into contiguous chunks, one per rank.
% Ranks are numbered clockwise from top-left (0,1,2,3), as in the original slide.
\documentclass[dvisvgm,border=3pt]{standalone}
\input{preamble.tex}
\begin{document}
\begin{tikzpicture}

  % --- root's send buffer, holding all 8 elements ---------------------------
  \mpibuf{(0,1.45)}{0,1,2,3,4,5,6,7}
  % curved arrow from the send buffer into root's own receive buffer
  \draw[flow] (-0.18,1.71) .. controls +(-0.75,0) and +(-0.75,0) .. (-0.18,0.96);

  % --- receive buffers: rank i receives elements 2i and 2i+1 ----------------
  \mpibuf{(0,0.70)}{0,1,.,.,.,.,.,.}          % rank 0 (top left, root)
  \mpibuf{(\colb,0.70)}{.,.,2,3,.,.,.,.}      % rank 1 (top right)
  \mpibuf{(\colb,-4.15)}{.,.,.,.,4,5,.,.}     % rank 2 (bottom right)
  \mpibuf{(0,-4.15)}{.,.,.,.,.,.,6,7}         % rank 3 (bottom left)

  % --- ranks and the scatter fan-out ---------------------------------------
  \mpiranks
  \mpiselfloop
  \mpiflows

\end{tikzpicture}
\end{document}
